\section{Poles with many paths}\label{sec:poles}

A cut of a cubic graph divides it into two parts, and a cycle that crosses the cut is made of
paths in the two parts. So we need to count paths between prescribed ends.

\begin{idea}{the results of this section}
First, if more than half of the possible paths between two ends are present, by a fixed margin,
then the part looks like a ladder, and it has few cycles (Theorem~\ref{thm:fewcycles}). Second,
the number of paths between two ends is the largest possible, or it is smaller by a fixed factor
(Theorem~\ref{thm:gap}).
\end{idea}

\subsection{Poles}

A \emph{pole} is a connected multigraph without loops in which every vertex has degree two or
three. Its vertices of degree two are its \emph{ports}. A \emph{$k$-pole} is a pole with exactly
$k$ ports. By a count of edge ends, a $k$-pole of rank $h$ has $2h+k-2$ vertices.

For two vertices $s\ne t$ of a multigraph $N$ let $P_{st}(N)$ be the number of paths from $s$ to
$t$. An \emph{$\{s,t\}$-join} is a set of edges in which $s$ and $t$ have odd degree and all other
vertices have even degree. If $N$ is connected of rank $b$, there are exactly $2^{b}$ of them: they
exist, and the sum of two of them is even. Every path from $s$ to $t$ is an $\{s,t\}$-join, so
\begin{equation}\label{eq:joinbound}
  P_{st}(N)\;\le\;2^{b}.
\end{equation}
Call the pair $s,t$ \emph{full} in $N$ if $P_{st}(N)=2^{\beta(N)}$.

\begin{lemma}[Joins]\label{lem:joins}
Let $N$ be a pole of rank $b$, and let $s,t$ be two of its ports. Every $\{s,t\}$-join is the union
of a path from $s$ to $t$ and of cycles, and these are pairwise vertex-disjoint. So the pair is
full if and only if every cycle of $N$ meets every path from $s$ to $t$. If some cycle of $N$ meets
no path from $s$ to $t$, then $P_{st}(N)\le2^{b-1}$.
\end{lemma}

\begin{proof}
In a join, $s$ and $t$ have odd degree at most two, and every other vertex has even degree at most
three. So the degrees are one at $s$ and $t$ and zero or two elsewhere. If a cycle $Z$ meets no
path from $s$ to $t$, then $J\mapsto J+Z$ maps the paths injectively to joins that are not paths.
\end{proof}

\begin{lemma}[Reduction]\label{lem:reduction}
Let $N$ be a pole of rank $b$, and let $s,t$ be two of its ports with $P_{st}(N)>2^{b-1}$.
\begin{enumerate}
\item[\textup{(a)}] Replace each port other than $s,t$, with its two edges, by one edge. No loop
arises. The result $N^{\circ}$ is a $2$-pole with ports $s,t$. It has the same rank, the same number
of paths from $s$ to $t$, and the same number of cycles as $N$.
\item[\textup{(b)}] Every bridge of $N^{\circ}$ separates $s$ from $t$ and is incident with neither.
Let $N_1,\dots,N_q$ be the components of $N^{\circ}$ without its bridges, in the order from $s$ to
$t$. Each $N_i$ is a bridgeless $2$-pole of rank $b_i\ge1$, and
\[
  b=\sum_ib_i,\qquad P_{st}(N)=\prod_iP(N_i),\qquad c(N)=\sum_ic(N_i),
\]
where $P(N_i)$ is the number of paths between the two ports of $N_i$. In particular
$P(N_i)\,2^{-b_i}\ge P_{st}(N)\,2^{-b}$ for every $i$.
\end{enumerate}
\end{lemma}

\begin{proof}
(a) Let $v$ be a port other than $s,t$. Suppose that both edges at $v$ end in the same vertex $u$.
Then $u$ has degree three, its third edge is a bridge, and no path from $s$ to $t$ meets the cycle
through $u$ and $v$. This contradicts Lemma~\ref{lem:joins}. So the two neighbours of $v$ are
distinct, and we may replace $v$ by an edge. Paths and cycles correspond. Repeat.

(b) Let $e$ be a bridge of $N^{\circ}$, and let $Y$ be a component of $N^{\circ}-e$, with $n$
vertices, $j$ of them ports. Counting edge ends gives $\beta(Y)=(n-j+1)/2$. This is a positive
integer, so $Y$ contains a cycle. If $Y$ contains neither port or both ports, a cycle on one side
of $e$ meets no path from $s$ to $t$. If $e$ is incident with $s$, every path from $s$ to $t$
starts with $e$ and avoids the cycle on the side of $s$. Both cases contradict
Lemma~\ref{lem:joins}. So every path from $s$ to $t$ passes all bridges, in the same order.

The multigraph $N_i$ is connected and bridgeless. Its vertices have degree three, except the ends
of the two bridges at $N_i$, or $s$, or $t$. These two vertices are distinct: a common end of two
bridges would have a third edge that is a bridge and does not separate $s$ from $t$. So $N_i$ is a
$2$-pole, and it contains a cycle. Ranks add over the $N_i$, every cycle lies in one $N_i$, and a
path from $s$ to $t$ is a concatenation of paths between the ports of the $N_i$ and of the bridges.
The last statement holds because $P(N_j)\le2^{b_j}$ for all $j$.
\end{proof}

\subsection{The shape of a pole with many paths}\label{sec:shape}

Let $N$ be a bridgeless $2$-pole with ports $s,t$ and rank $b\ge2$. It has $2b$ vertices. It has no
cut vertex: a cut vertex of a subcubic multigraph is joined to some component by a single edge.
We use the following notation. The \emph{layer} $L_j$ is the set of vertices at distance $j$ from
$s$. Every edge joins vertices of the same layer or of consecutive layers. Let $\ell$ be the
distance from $s$ to $t$, let $A=w_0w_1\cdots w_\ell$ be a shortest path from $s$ to $t$, and put
$J=N-V(A)$. Then $w_j\in L_j$. A vertex of $A$ lies on two edges of $A$, or on one if it is a
port. So it has at most one \emph{further edge}. This edge leads to $J$, or it is parallel to an
edge of $A$, because $A$ is a shortest path. For $0<j<\ell$ the layer $L_j$ separates $s$ from $t$,
so $|L_j|\ge2$. Call $L_j$ \emph{thin}
if $0<j<\ell$ and $|L_j|=2$. A thin layer consists of $w_j$ and one vertex of $J$. Call a thin
layer \emph{ordinary} if its vertex of $J$ is on no pair of parallel edges
(Figure~\ref{fig:layers}).

\begin{figure}[ht]
\centering
\begin{tikzpicture}[x=1.22cm,y=1cm]
  % layers
  \foreach \j in {1,2,3,6,7,8}
    \draw[slate!45, line width=0.4pt, dash pattern=on 2pt off 2pt] (\j,-0.45) -- (\j,1.75);
  \foreach \j in {1,2,3,6,7,8} \node[note] at (\j,2.0) {$L_{\j}$};
  \node[note] at (0,2.0) {$L_0$};
  \node[note] at (9,2.0) {$L_\ell$};
  % the part of bounded size
  \draw[draw=warm, fill=warm!8, line width=0.6pt, rounded corners=7pt] (3.55,-0.4) rectangle (5.45,1.7);
  \node[note, text=warm, align=center] at (4.5,0.85) {not thin:\\ bounded size};
  % shortest path A
  \foreach \j in {0,...,9} \coordinate (w\j) at (\j,0);
  \foreach \j in {1,2,3,6,7,8} \coordinate (y\j) at (\j,1.3);
  \draw[cool] (w0)--(w1)--(w2)--(w3)--(3.55,0) (5.45,0)--(w6)--(w7)--(w8)--(w9);
  \draw[cool, dash pattern=on 2.5pt off 2pt] (3.55,0)--(5.45,0);
  % edges of J and between
  \draw[edge] (w0)--(y1) (y1)--(y2) (w1)--(y2) (y1)--(w2) (y2)--(y3) (w3)--(y3) (y3)--(3.55,1.3);
  \draw[edge] (5.45,1.3)--(y6) (w6)--(y6) (y6)--(y7) (y7)--(y8) (w7)--(y8) (y7)--(w8) (y8)--(w9);
  \foreach \j in {1,2,3,6,7,8} \node[vtx] at (w\j) {};
  \foreach \j in {1,2,3,6,7,8} \node[vtx] at (y\j) {};
  \node[port] at (w0) {};
  \node[port] at (w9) {};
  \node[note, text=accent, below=3pt] at (w0) {$s$};
  \node[note, text=accent, below=3pt] at (w9) {$t$};
  \node[note, text=accent, below=2pt] at (2,0) {$w_2$};
  \node[note, text=accent, anchor=west] at (9.25,0) {the shortest path $A$};
  \node[note, anchor=west] at (9.25,1.3) {the rest $J$};
\end{tikzpicture}
\caption{A bridgeless $2$-pole with many paths from $s$ to $t$ (schematic). Almost every layer
is thin: it has one vertex of the shortest path $A$ and one vertex of $J$
(Proposition~\ref{prop:thin}).}
\label{fig:layers}
\end{figure}

\begin{idea}{outline of the proof of Proposition~\ref{prop:thin}}
Add the shortest path $A$ to a path from $s$ to $t$. The result is an even subgraph whose edges
away from $A$ form a forest. So more than half of the outside parts are forests, and every layer
is small by Lemma~\ref{lem:width}. By
Lemma~\ref{lem:smallcuts} there are few short cycles away from $A$. By a count of vertices,
almost every layer then has exactly two vertices.
\end{idea}

\begin{proposition}[Thin layers]\label{prop:thin}
For every $\varepsilon\in(0,\half]$ there is a number $C=C_\varepsilon\ge1$ with the following
property. Let $N$ be a bridgeless $2$-pole of rank $b\ge2$ with ports $s,t$ and
$P_{st}(N)\ge(\half+\varepsilon)2^{b}$. Then, in the notation above,
\begin{enumerate}
\item[\textup{(i)}] $\beta(J)\le C$ and $b-C\le\ell\le b$;
\item[\textup{(ii)}] at most $2C$ of the layers $L_j$ with $0<j<\ell$ are not thin, and
$L_j=\emptyset$ for $j>b+C$;
\item[\textup{(iii)}] at most $10C+2$ vertices of $N$ lie outside the ordinary layers.
\end{enumerate}
\end{proposition}

\begin{proof}
\step{Step 1: the layers are small.} Put $P=V(A)$. Then $N[P]$ is connected, and every vertex
outside $P$ has degree three. For a path $Q$ from $s$ to $t$, the set $Q+A$ is an even subgraph of
$N$. Its edges in $N-P$ are edges of the path $Q$, so they form a forest. Different paths give
different even subgraphs, and $N$ has $2^{b}$ even subgraphs. Hence
\begin{equation}\label{eq:qH}
  q(N,P)\;\ge\;\frac{P_{st}(N)}{2^{b}}\;\ge\;\frac12+\varepsilon .
\end{equation}
Let $\lambda(u)$ be the distance from $s$ to $u$ in $N$. This is a level function, and $P$ has at
most one vertex of each level. By Lemma~\ref{lem:width} with $\delta=\varepsilon/2$, every layer
$L_j$ has fewer than $\Theta:=\Theta_{\varepsilon/2}(1)$ vertices of $J$.

\step{Step 2: a short cycle.} Put $U_j=L_j\sm V(A)$, and for an interval $I$ of indices let $J_I$ be
the subgraph of $J$ induced by the sets $U_j$ with $j\in I$. Suppose that $J_I$ contains a cycle.
Choose a cycle $Z$ of $J_I$ such that the interval $[a,a']$ of the layers that $Z$ meets is as
short as possible, and put $\sigma=a'-a$ and $R=U_a\cup\dots\cup U_{a'}$. We claim that
$|R|\le20\Theta$.

By the choice of $Z$, every cycle of $J[R]$ meets $U_a$. So $J[R]-U_a$ is a forest. Deleting a
vertex of degree at most three lowers the rank by at most two, so $\beta(J[R])\le2|U_a|<2\Theta$.
The vertices of $R$ have degree three in $N$. Counting edge ends gives
$|R|=2\beta(J[R])+|\delta R|-2\comp(J[R])$. An edge that leaves $R$ ends in $U_{a-1}$, in
$U_{a'+1}$, or in one of the vertices $w_{a-1},\dots,w_{a'+1}$, each of which has at most one
further edge. So $|\delta R|<6\Theta+\sigma+3$, and $|R|<10\Theta+\sigma+1$. The cycle $Z$ has a vertex in
layer $a$ and a vertex in layer $a'$, so $|Z|\ge2\sigma$. Also $|Z|\le|R|$. Hence
$\sigma\le10\Theta$ and $|R|\le20\Theta$.

\step{Step 3: the rank of $J$.} Start with the interval of all layers. As long as $J_I$ contains a
cycle for some interval $I$ of the present family, choose $Z$ and $R$ as in Step~2, and replace $I$
by the at most two intervals that remain when $[a,a']$ is removed. Let $Z_1,\dots,Z_k$ be the
chosen cycles and $R_1,\dots,R_k$ the corresponding sets. The sets $R_i$ are pairwise disjoint,
$|R_i|\le20\Theta$, and $J-\bigcup_iR_i$ is a forest, because an edge of $J$ joins equal or
consecutive layers. So $\beta(J)\le2\sum_i|R_i|\le40\Theta k$.

The sets $V(Z_i)$ are pairwise disjoint cyclic regions of $N$ that avoid $P$. Each has at most
$|Z_i|\le20\Theta$ edges that leave it. If $\xi\cap E(N-P)$ is a forest, then
$\xi\cap E(N[V(Z_i)])$ is a forest for every $i$. So Lemma~\ref{lem:smallcuts} and~\eqref{eq:qH}
give $k\le(20\Theta+1)2^{20\Theta-1}/\varepsilon$. Hence $\beta(J)\le C$ with
\[
  C=C_\varepsilon:=\frac{40\Theta\,(20\Theta+1)\,2^{20\Theta-1}}{\varepsilon}.
\]

\step{Step 4: counting.} The layers $L_0,\dots,L_\ell$ contain at least $2\ell$ vertices, so
$\ell\le b$. The vertices of $A$ meet the $\ell$ edges of $A$ and at most $\ell+1$ other edges.
If $J$ is not empty, then $|E(J)|\ge(3b-1)-(2\ell+1)$ and $|V(J)|=2b-\ell-1$, so
$\beta(J)\ge|E(J)|-|V(J)|+1\ge b-\ell$. If $J$ is empty, then $2b=\ell+1$. In both cases
$\Delta:=b-\ell\le C$. This is~(i).

Compare $N$ with a multigraph that has one vertex in $L_0$ and in $L_\ell$, two in each layer
between, and none beyond. Then $N$ has exactly $2\Delta$ vertices more. So at most $2\Delta$
layers between $L_0$ and $L_\ell$ are not thin, and at most $2\Delta$ layers lie beyond $L_\ell$.
This is~(ii). A layer between $L_0$ and $L_\ell$ with $w\ge3$ vertices has $w-2\ge w/3$ vertices
too many. So at most $2+3\cdot2\Delta$ vertices lie outside the thin layers.

Finally let $y\in J$ be on a pair of parallel edges. The other end of the pair has two edges to
$y$, so it is not on $A$, whose vertices have at most one further edge. So the pair is a cycle of
$J$. Two such pairs have no common vertex,
because $N$ is subcubic and has more than two vertices. So there are at most $\beta(J)\le C$ pairs,
and at most $2C$ thin layers are not ordinary. They contain at most $4C$ vertices. This gives~(iii).
\end{proof}

\begin{theorem}[Few cycles]\label{thm:fewcycles}
For every $\varepsilon\in(0,\half]$ there is a constant $K_\varepsilon$ with the following property.
Let $N$ be a pole of rank $b$, and let $s,t$ be two of its ports with
$P_{st}(N)\ge(\half+\varepsilon)2^{b}$. Then
\[
  c(N)\;\le\;K_\varepsilon\,2^{b/2}.
\]
\end{theorem}

\begin{proof}
Let $C=C_\varepsilon$ be as in Proposition~\ref{prop:thin}. We prove that a bridgeless $2$-pole
with these properties has at most $2^{31C+22}\,2^{b/2}$ cycles. The theorem follows with
$K_\varepsilon=2^{31C+23}$. Indeed, in the notation of Lemma~\ref{lem:reduction}, each $N_i$ has the
required number of paths, and $\sum_i2^{b_i/2}\le2^{b/2}+q-1\le2\cdot2^{b/2}$, because
$x+y\le xy+1$ for $x,y\ge1$.

So let $N$ be a bridgeless $2$-pole. If $b=1$, then $N$ has one cycle. Let $b\ge2$, and use the
notation above. Let $Z$ be a cycle, and let $L_i$ and $L_j$ be the first and the last layer that it
meets. Let $M$ be the set of the vertices that lie in $L_i$, in $L_j$, or outside the ordinary
layers. By Proposition~\ref{prop:thin}(iii), $|M|\le10C+6$.

\step{Two strands.} Let $L_l=\{w_l,y_l\}$ be an ordinary layer with $i<l<j$. The cycle $Z$ has
vertices before and after this layer. So it has at least two edges from $L_{l-1}$ to $L_l$ and at
least two edges from $L_l$ to $L_{l+1}$. The two vertices of $L_l$ have degree at most two in $Z$.
So $Z$ has exactly two edges on each side of $L_l$ and does not contain an edge $w_ly_l$. Suppose
that both edges from $L_{l-1}$ end in the same vertex of $L_l$. Then both edges to $L_{l+1}$ start
in the other vertex, and no edge of $Z$ joins the part of $Z$ before $L_l$ to the part after it.
This is impossible, because $Z$ is connected. So each vertex of $L_l$ has one edge of $Z$ to
$L_{l-1}$ and one to $L_{l+1}$.

\step{Matchings.} Let $L_l$ and $L_{l+1}$ be ordinary layers with $i<l$ and $l+1<j$. By the last
paragraph, the edges of $Z$ between them form a perfect matching between $\{w_l,y_l\}$ and
$\{w_{l+1},y_{l+1}\}$. We count the perfect matchings. The edge $y_ly_{l+1}$ is not doubled,
because the layers are ordinary. The edges $w_ly_{l+1}$ and $y_lw_{l+1}$ are not doubled, because
a vertex of $A$ has at most one further edge. If the edge $w_lw_{l+1}$ of $A$ is doubled, then
$w_l$ and $w_{l+1}$ have no edge to $J$. So there are at most two perfect matchings. If there
are two, then $w_{l+1}$ has two edge ends towards $L_l$: either both $w_l$ and $y_l$ are its
neighbours, or $w_lw_{l+1}$ is doubled. In the same way $w_{l+1}$ has two edge ends towards
$L_{l+2}$ if there are two perfect matchings between $L_{l+1}$ and $L_{l+2}$. Both cannot happen,
because $w_{l+1}$ has degree three. So among consecutive pairs of ordinary layers, no two
neighbouring pairs have two perfect matchings each.

\step{The count.} An edge of $Z$ has an end in $M$, or it joins two consecutive ordinary layers
strictly between $L_i$ and $L_j$. So $Z$ is determined by $i$ and $j$, by its edges at $M$, and by
its perfect matchings. There are at most $2^{3|M|}$ sets of edges at $M$. By the last paragraph
there are at most $2^{(j-i)/2}$ choices for the matchings. All layers have an index of at most
$n:=b+C$, by Proposition~\ref{prop:thin}(ii). Hence
\[
  c(N)\;\le\;2^{30C+18}\sum_{0\le i\le j\le n}2^{(j-i)/2}
  \;\le\;2^{30C+18}\cdot12\cdot2^{(b+C)/2}\;\le\;2^{31C+22}\,2^{b/2}.\qedhere
\]
\end{proof}

A pole with many paths also has a short cycle:

\begin{lemma}[Four thin layers]\label{lem:fourlayers}
Let $N$ be a bridgeless $2$-pole of rank at least two. If four consecutive layers are thin, then
$N$ has a cycle of length at most four all of whose vertices lie in these four layers.
\end{lemma}

\begin{proof}
The eight vertices of the four layers have degree three, and each has a neighbour in the layer
before its own. The edges between two consecutive layers form a cut of $N$, so there are at least
two of them. There are at most four: the two vertices of the earlier layer have six edge ends, and
two of these point backwards. A pair of parallel edges among the eight vertices is a cycle of
length two. Suppose there is none. If some cut between two of the four layers has four edges, it
is a cycle of length four. So let the three cuts have $e_1,e_2,e_3\in\{2,3\}$ edges. Let $r_2$ and
$r_3$ be the numbers of edges inside the second and the third layer. Counting edge ends in these
layers gives $e_1+e_2+2r_2=6$ and $e_2+e_3+2r_3=6$. So $e_1=e_2=e_3$.

If all three are $2$, then $r_2=r_3=1$, and each vertex of the second and third layer has one edge
to each side. The two edges between the second and the third layer and the edges inside these
layers form a cycle of length four. If all three are $3$, then $r_2=0$. In the second layer one
vertex $u$ has two neighbours in the first layer and the other vertex $u'$ has one, say $x$; then
$x$ is also a neighbour of $u$. So $u$ has one neighbour $y$ in the third layer and $u'$ has two,
among them $y$. Then $xuyu'$ is a cycle of length four.
\end{proof}

\begin{corollary}\label{cor:nodense}
For every $\varepsilon\in(0,\half]$ and every integer $g\ge0$ there is $b_1(\varepsilon,g)$ with the
following property. Let $G$ be a connected cubic multigraph without loops, of rank at least $9$,
that has no cyclic cut of size at most four. Let $K$ be an induced subgraph of $G$ that is a pole
with at most $g+2$ ports and rank $\beta(K)\ge b_1(\varepsilon,g)$. Then
$P_{st}(K)<(\half+\varepsilon)2^{\beta(K)}$ for all ports $s,t$ of $K$.
\end{corollary}

\begin{proof}
Suppose that $P_{st}(K)\ge(\half+\varepsilon)2^{\beta(K)}$, and let $N=K^{\circ}$ and
$N_1,\dots,N_q$ be as in Lemma~\ref{lem:reduction}. An edge of $N$ is a path of $K$ whose inner
vertices are ports of $K$. Call the edge \emph{plain} if this path has no inner vertex. At most
$g$ edges of $N$ are not plain.

Let $1<i<q$, and suppose that all edges of $N_i$ are plain. Then the vertices of $N_i$ have degree
three in $K$. So exactly two edges of $G$ leave $V(N_i)$, one for each bridge at $N_i$. The
multigraph $N_i$ contains a cycle, and so does $N_1$. This is a cyclic cut of size two, which $G$
does not have. So $q\le g+2$, and some $N_i$ has rank $b'\ge\beta(K)/(g+2)$.

Let $C=C_\varepsilon$. By Proposition~\ref{prop:thin}, applied to $N_i$, the distance between its
ports is at least $b'-C$, and at most $2C$ of its inner layers are not thin. At most $2g$ layers
contain an end of an edge that is not plain. Divide the inner layers into groups of four
consecutive layers. If $b'\ge9C+8g+9$, then in some group all layers are thin and all edges with
both ends in the group are plain. Lemma~\ref{lem:fourlayers} gives a cycle $Z$ of length at most
four in this group. It is a cycle of $G$. At most four edges of $G$ leave $V(Z)$, and at most
eight edges of $G$ have an end in $V(Z)$. So $G-V(Z)$ has rank at least $\beta(G)-8>0$, and
$\delta V(Z)$ is a cyclic cut of size at most four. This is a contradiction. So the corollary
holds with $b_1(\varepsilon,g)=(g+2)(9C+8g+9)$.
\end{proof}

\paragraph{Consequences of Theorem~\ref{thm:fewcycles}.}
From now on fix
\[
  \varepsilon_0=\tfrac1{1000},\qquad K_0=K_{\varepsilon_0},\qquad\kappa(h)=2K_0\,2^{-h/2}.
\]
For four distinct ports $a,b,c,d$ of a pole $N$ let $f_{ab|cd}(N)$ be the number of pairs of
vertex-disjoint paths of $N$, one from $a$ to $b$ and one from $c$ to $d$.

\begin{lemma}\label{lem:gadgets}
Let $N$ be a pole of rank $h$. In each statement below, the letters denote distinct ports of $N$.
Write $P_{xy}$ for $P_{xy}(N)$ and $f$ for $f(N)$.
\begin{enumerate}
\item[\textup{(a)}] If $P_{ab}\ge0.501\cdot2^{h}$, then $c(N)\le K_0\,2^{h/2}$.
\item[\textup{(b)}] If $P_{ac}+P_{bc}\ge1.002\cdot2^{h}$, then $P_{ab}\le\kappa(h)\,2^{h}$.
\item[\textup{(c)}] If $P_{ac}+P_{ad}+P_{bc}+P_{bd}\ge2.004\cdot2^{h}$, then
$f_{ac|bd}+f_{ad|bc}\le\kappa(h)\,2^{h}$.
\item[\textup{(d)}] If $P_{ab}+f_{ac|bd}+f_{ad|bc}\ge1.002\cdot2^{h}$, then
$P_{cd}\le\kappa(h)\,2^{h}$.
\end{enumerate}
\end{lemma}

\begin{proof}
(a) is Theorem~\ref{thm:fewcycles}. In each of the other cases we change $N$ into a pole $N'$ and
apply Theorem~\ref{thm:fewcycles} to two ports of $N'$.

(b) Add a new vertex $z$ and join it to $a$ and $b$. Then $N'$ has rank $h+1$, and
$P_{zc}(N')=P_{ac}+P_{bc}\ge(\half+\varepsilon_0)2^{h+1}$. The cycles of $N'$ through $z$
correspond to the paths of $N$ from $a$ to $b$. So $P_{ab}\le c(N')\le K_0\,2^{(h+1)/2}$.

(c) Add a new vertex $z$ joined to $a$ and $b$, and a new vertex $w$ joined to $c$ and $d$. Then
$N'$ has rank $h+2$, and $P_{zw}(N')$ is the sum of the four path counts. The cycles of $N'$
through $z$ and $w$ correspond to the pairs counted by $f_{ac|bd}+f_{ad|bc}$. So this sum is at
most $K_0\,2^{(h+2)/2}$.

(d) Add an edge between $c$ and $d$. Then $N'$ has rank $h+1$. A path from $a$ to $b$ in $N'$
avoids the new edge, or it consists of the new edge and a pair counted by
$f_{ac|bd}+f_{ad|bc}$. So $P_{ab}(N')=P_{ab}+f_{ac|bd}+f_{ad|bc}$. The cycles of $N'$ through the
new edge correspond to the paths of $N$ from $c$ to $d$. So $P_{cd}\le K_0\,2^{(h+1)/2}$.
\end{proof}

\subsection{Full pairs}\label{sec:full}

A \emph{square} is a copy of the complete bipartite graph $K_{2,2}$; its two classes are its
\emph{left pair} and its \emph{right pair}. A \emph{single rung} is one edge; its two ends form
both its left pair and its right pair. These are the two kinds of \emph{tiles}. The \emph{weight}
of a square is $2$, and the weight of a single rung is $1$. A \emph{chain} of tiles
$\tau_1,\dots,\tau_m$ consists of these tiles, pairwise disjoint, and of two disjoint edges from
the right pair of $\tau_i$ to the left pair of $\tau_{i+1}$ for $1\le i<m$. These two edges are an
\emph{interface}. A \emph{capped ladder} is a chain with two more vertices $s$ and $t$, the
\emph{caps}: $s$ is joined to the two vertices of the left pair of $\tau_1$, and $t$ to the two
vertices of the right pair of $\tau_m$ (Figure~\ref{fig:ladder}). A capped ladder is a $2$-pole. If
its tiles have total weight $W$, it has $2W+2$ vertices and rank $W+1$.

\begin{figure}[ht]
\centering
\begin{tikzpicture}[x=1cm,y=1cm]
  \coordinate (s) at (0,0);
  \coordinate (t) at (9.6,0);
  \foreach \n/\x in {a/1.2,b/2.4,c/3.8,d/5.2,e/6.4,f/7.8}{
    \coordinate (\n u) at (\x,0.65);
    \coordinate (\n l) at (\x,-0.65);}
  \draw[edge] (s)--(au) (s)--(al);
  \draw[cool] (au)--(bu) (al)--(bl) (au)--(bl) (al)--(bu);
  \draw[edge] (bu)--(cu) (bl)--(cl);
  \draw[hot] (cu)--(cl);
  \draw[edge] (cu)--(du) (cl)--(dl);
  \draw[cool] (du)--(eu) (dl)--(el) (du)--(el) (dl)--(eu);
  \draw[edge] (eu)--(fu) (el)--(fl);
  \draw[hot] (fu)--(fl);
  \draw[edge] (fu)--(t) (fl)--(t);
  \foreach \q in {au,al,bu,bl,cu,cl,du,dl,eu,el,fu,fl} \node[vtx] at (\q) {};
  \node[port] at (s) {};
  \node[port] at (t) {};
  \node[note, text=accent, left=3pt] at (s) {$s$};
  \node[note, text=accent, right=3pt] at (t) {$t$};
  \node[note, text=accent] at (1.8,-1.15) {square};
  \node[note, text=warm] at (3.8,-1.15) {single rung};
  \node[note, text=accent] at (5.8,-1.15) {square};
  \node[note, text=warm] at (7.8,-1.15) {single rung};
  \node[note] at (3.1,1.1) {interface};
  \draw[slate, line width=0.4pt] (3.1,0.95) -- (3.1,0.7);
\end{tikzpicture}
\caption{A capped ladder with four tiles, of weight $6$ and rank $7$. Without $s$ and $t$ it is a
chain. If the right pair of the last tile is joined to the left pair of the first tile instead, the
result is a ring (Section~\ref{sec:rings}).}
\label{fig:ladder}
\end{figure}

\begin{lemma}[Ladders]\label{lem:ladders}
Let $N$ be a bridgeless $2$-pole whose pair of ports is full. Then $N$ consists of two vertices and
two parallel edges, or $N$ is a capped ladder whose caps are the ports.
\end{lemma}

\begin{proof}
Let $s,t$ be the ports. By Lemma~\ref{lem:joins}, every cycle meets every path from $s$ to $t$. If
the rank is one, $N$ is a cycle with two vertices. Otherwise $N$ has no cut vertex, so there are
two paths $Q$ and $Q'$ from $s$ to $t$ without a common inner vertex. Let $F$ be the set of the
edges outside $Q\cup Q'$. In $F$ the inner vertices of $Q$ and $Q'$ have degree one, the other
vertices except $s,t$ have degree three, and $s,t$ have degree zero.

A cycle in $F$ would avoid $Q$. So $F$ is a forest. Suppose that a component of $F$ has two leaves
on $Q$. The path of $F$ between them and a part of $Q$ form a cycle that avoids $Q'$. This is
impossible, and the same holds for $Q'$. So every component of $F$ has at most two leaves. A tree
with degrees one and three has two more leaves than vertices of degree three. So every component
of $F$ with an edge is a single edge, from an inner vertex of $Q$ to an inner vertex of $Q'$. Hence
all vertices lie on $Q\cup Q'$, and $F$ is a perfect matching between the inner vertices
$x_1,\dots,x_k$ of $Q$ and $y_1,\dots,y_k$ of $Q'$, in the order from $s$ to $t$. Write its edges
as $x_iy_{\pi(i)}$.

Let $i<j<l$. If $\pi(l)<\min\{\pi(i),\pi(j)\}$, then the matching edges at $x_i$ and $x_j$ and the
parts of $Q$ and $Q'$ between their ends form a cycle. The path that follows $Q'$ from $s$ to
$y_{\pi(l)}$, crosses to $x_l$ and follows $Q$ to $t$ avoids it. If
$\pi(i)>\max\{\pi(j),\pi(l)\}$, the same holds for the cycle through the matching edges at $x_j$
and $x_l$ and the path that follows $Q$ to $x_i$, crosses, and follows $Q'$ to $t$. So neither
case occurs. It follows that $\pi(1)\le2$, and that $\pi(2)=1$ if $\pi(1)=2$. By induction, $\pi$
is composed of fixed points and of exchanges of two consecutive numbers. A fixed point $i$ is a
single rung with the ends $x_i,y_i$. An exchange of $i$ and $i+1$ is a square with the left pair
$x_i,y_i$ and the right pair $x_{i+1},y_{i+1}$. The remaining edges of $Q$ and $Q'$ are the
interfaces and the edges at the caps.
\end{proof}

We need two counts in a capped ladder. Let $\tau_1,\dots,\tau_m$ be its tiles. First, let
$1\le i\le m$ and let $u$ be a vertex of the left pair of $\tau_i$. The number of paths from $u$ to
$t$ inside $\tau_i\cup\dots\cup\tau_m\cup\{t\}$ is $2^{q}$, where $q$ is the weight of
$\tau_i,\dots,\tau_m$. Indeed, such a path passes the tiles in order. It enters a tile at a vertex
of the left pair and leaves it at a vertex of the right pair. In a square there are two paths to
each vertex of the right pair. In a single rung the path stays or crosses. Second, let $A_i$ be
the number of paths between the two vertices of the right pair of $\tau_i$ inside
$\{s\}\cup\tau_1\cup\dots\cup\tau_i$, and $A_0=1$. This multigraph has rank equal to the weight
$p$ of $\tau_1,\dots,\tau_i$. So $A_i\le2^{p}$ by~\eqref{eq:joinbound}.

\begin{lemma}[A third port]\label{lem:thirdport}
Let $T$ be a $3$-pole of rank $h$ with ports $s,t,v$. If the pair $s,t$ is full, then
$P_{vt}(T)=2^{h}$ or $P_{vt}(T)\le\frac78\,2^{h}$.
\end{lemma}

\begin{proof}
By Lemma~\ref{lem:reduction}, $T$ arises from a $2$-pole $N$ with ports $s,t$ when an edge $e$ is
subdivided by $v$, and $N$ consists of bridges and of bridgeless $2$-poles $N_1,\dots,N_q$ with
full pairs. If $e$ is a bridge or an edge of some $N_i$ with $i\ge2$, then no path from $v$ to $t$
meets $N_1$. Such a path lies in a part of $T$ of rank at most $h-1$, so $P_{vt}(T)\le2^{h-1}$
by~\eqref{eq:joinbound}. Let $e$ be an edge of $N_1$, and let $t_1$ be the second port of $N_1$. A
path from $v$ to $t$ is a path from $v$ to $t_1$ followed by paths through $N_2,\dots,N_q$, and
each of these $2$-poles is full. So we may assume that $q=1$.

By Lemma~\ref{lem:ladders}, $N$ is a cycle with two vertices or a capped ladder. In the first
case $T$ is a triangle and $P_{vt}=2=2^{h}$. In the second case let $p$ be the weight of the
tiles before $e$ and $q'$ the weight of the tiles after $e$; if $e$ lies in a tile, this tile is
not counted. A path from $v$ to $t$ never returns once it has entered the tiles after $e$, and
from its point of entry it has $2^{q'}$ continuations. There are three cases.

\step{The edge $e$ is in an interface or at a cap.} Then $h=p+q'+1$. The path goes from $v$ to
the right, or it goes to the left, follows one of the $A$ paths between the two ends of the
interface on the left side, and uses the other edge of the interface. Here $A=A_i$ for the tile
$\tau_i$ before $e$. So $P_{vt}=(1+A)2^{q'}$. If $p=0$, this is $2^{h}$. If $p\ge1$, then
$A\le2^{p}$ and $P_{vt}\le\frac34\,2^{h}$.

\step{The edge $e$ is a single rung.} Then $h=p+q'+2$. From each end of the rung the path
continues to the right, or it makes the detour through the tiles before the rung to the other end
of the rung and continues from there. So $P_{vt}=2(1+A)2^{q'}$, and the conclusion is the same.

\step{The edge $e$ is in a square.} Then $h=p+q'+3$. Let $u,u'$ be the left pair and $x,x'$ the
right pair of the square, and let $e=ux$. A path from $v$ leaves the square at $x$ or $x'$.
Without a detour through the tiles before the square, its part in the square is one of
\[
  vx,\qquad vux'u'x,\qquad vux',\qquad vxu'x' .
\]
With a detour between $u$ and $u'$ it is one of $vu\cdots u'x$, $vu\cdots u'x'$ and
$vxu'\cdots ux'$. There are $A$ detours. So $P_{vt}=(4+3A)2^{q'}$. If $p=0$, then $A=1$ and
$P_{vt}=\frac78\,2^{h}$. If $p\ge1$, then $P_{vt}\le(4+3\cdot2^{p})2^{q'}\le\frac58\,2^{h}$.
\end{proof}

\begin{theorem}[Gap]\label{thm:gap}
Let $N$ be a pole of rank $b$, and let $s,t$ be two of its ports. Then
\[
  P_{st}(N)=2^{b}\qquad\text{or}\qquad P_{st}(N)\le\tfrac{15}{16}\,2^{b}.
\]
\end{theorem}

\begin{proof}
We use induction on $b$. Let $P_{st}(N)>2^{b-1}$. By Lemma~\ref{lem:reduction} we may assume that
$N$ is a $2$-pole. If $N$ has a bridge, the claim holds for the multigraphs $N_i$ of
Lemma~\ref{lem:reduction}, which have smaller rank, and the path counts multiply. So let $N$ be
bridgeless. If $b=1$, the pair is full. Let $b\ge2$. If $s$ and $t$ are adjacent, then
$P_{st}(N)\le1+2^{b-1}\le\frac34\,2^{b}$, by~\eqref{eq:joinbound} for $N$ without that edge.
Otherwise let $x$ and $y$ be the neighbours of $s$. They are distinct, because $N$ has no bridge.
The multigraph $T=N-s$ is connected, because $N$ has no cut vertex. It is a $3$-pole of rank $b-1$
with ports $x,y,t$, and
\[
  P_{st}(N)=P_{xt}(T)+P_{yt}(T).
\]
By induction each term is $2^{b-1}$ or at most $\frac{15}{16}2^{b-1}$. If neither pair is full, we
are done. If the pair $x,t$ is full, then $P_{yt}(T)$ is $2^{b-1}$ or at most $\frac78\,2^{b-1}$,
by Lemma~\ref{lem:thirdport}. So $P_{st}(N)$ is $2^{b}$ or at most $\frac{15}{16}2^{b}$.
\end{proof}

\begin{corollary}[Four ports]\label{cor:fourports}
Let $N$ be a $4$-pole with ports $a,b,c,d$ such that the four pairs $ac$, $ad$, $bc$, $bd$ are
full. Then one of the following holds.
\begin{enumerate}
\item[\textup{(i)}] $N$ is a chain of tiles; $a,b$ is the left pair of its first tile and $c,d$ is
the right pair of its last tile.
\item[\textup{(ii)}] $N$ has a bridge $e$ such that one component of $N-e$ contains $a$, $b$ and a
cycle, and the other component contains $c$, $d$ and a cycle.
\end{enumerate}
\end{corollary}

\begin{proof}
Add a vertex $z$ joined to $a,b$ and a vertex $w$ joined to $c,d$. The result $N'$ is a $2$-pole
of rank $\beta(N)+2$ with ports $z,w$, and $P_{zw}(N')=4\cdot2^{\beta(N)}$. So the pair $z,w$ is
full. Let $N_1,\dots,N_q$ be as in Lemma~\ref{lem:reduction}. Each $N_i$ has a full pair of ports,
so Lemma~\ref{lem:ladders} applies to it. The neighbours $a,b$ of $z$ are distinct, so $N_1$ is a
capped ladder with the cap $z$, and $a,b$ is the left pair of its first tile. In the same way
$N_q$ is a capped ladder with the cap $w$. If $q=1$, this is~(i). If $q\ge2$, let $e$ be the
bridge at $N_1$. The multigraphs $N_1-z$ and $N_q-w$ are connected, each contains a cycle, and they
lie in different components of $N-e$. This is~(ii).
\end{proof}
